% ===================================================================
%  図表第 5 号 — 家とその建て方。
%
%  Ceci n'est PAS une transcription : c'est une traduction, faite en
%  2026, en japonais d'aujourd'hui. Voir texto/ja/10-zuhyo-01.tex
%  pour la consigne, qui vaut pour les seize fichiers.
%
%  LES DEUX NOMS SONT TRANSCRITS, non traduits : le japonais rend les
%  noms francais en katakana, et il l'a fait pour tous les prenoms de
%  ces seize tableaux. M. Leblanc est « ルブラン氏 », M. Roux est
%  « ルー氏 ». Le chinois, lui, traduisait la couleur — 白先生 — parce
%  qu'il n'a pas de syllabaire pour transcrire.
% ===================================================================

%%K t05-apar-1 apar
\VUsaut{6.0mm}
\VUtitre{15.0pt}{60}{図表第 5 号}
\VUsaut{1.4mm}
\VUtitre{11.4pt}{0}{家とその建て方。}
\VUsaut{2.2mm}

%%K t05-tit-1 sub
\VUpk{10.2pt}{40}{よい出会い。}
\VUsaut{3.0mm}

%%K t05-01-1 p
\textsc{ジャン}。--- 叔父さん、\VUgras{家}がどう建つのかとても知りたいのです。

%%K t05-01-2 p
\textsc{叔父} \textsuperscript{(80)}。--- それはたやすい、坊や。ついておいで。
ベルナール氏 \textsuperscript{(79)} が建てさせている家を見せてやろう。私がこれから
住む家だ。

%%K t05-01-3 p
\textsc{ジャン}。--- ではこの家の\VUgras{図面} \textsuperscript{(76)} は誰が描いた
のですか。

%%K t05-01-4 p
\textsc{叔父}。--- \VUgras{建築家} \textsuperscript{(75)} だ。たいそう明らかな図を
出して、それが通り、\VUgras{請負人} \textsuperscript{(77)} が工事を始めた。

%%K t05-01-5 p
\textsc{ジャン}。--- 請負人とは誰ですか。

%%K t05-01-6 p
\textsc{叔父}。--- ルブラン氏だ、お前の友だちジャックの父上だよ。建築家の
ルー氏と二人で職人を指図しておられる。

%%K t05-01-7 p
おや、ちょうどルブラン氏が自分の\VUgras{物差} \textsuperscript{(78)} を持って
来られる。ルー氏も自分の図面を持って来られる。

%%K t05-01-8 p
おはようございます、御両所。いかがお過ごしですか。

%%K t05-01-9 p
\textsc{ルー氏}。--- たいそう結構です、ありがとう。おはようジャン。この子は
なんと大きくなったことか。

%%K t05-01-10 p
\textsc{叔父}。--- まったく、すっかり小さな大人です。たいそう真面目で、見る
ものは何でも説いてやらねばなりません。今日は家がどう建つのかを知りたいそう
です。

%%K t05-tit-2 sub
\VUpk{10.2pt}{40}{職人。}
\VUsaut{3.0mm}

%%K t05-01-11 p
\textsc{ルブラン氏}。--- ではついておいで、小さな友だち。すぐに説いてあげよう。
学びたがる子供が私はたいそう好きだからね。

%%K t05-01-12 p
あの手に\VUgras{鋤} \textsuperscript{(82)} を持った職人が見えるかね。あれは
\VUgras{土工} \textsuperscript{(81)} だ。水道の管を埋めるために\VUgras{溝}
\textsuperscript{(46)} を掘っている。家の建つ地面を掘ったのもあの男で、左官が四つ
の\VUgras{壁}を立てられるようにするためだ。その壁の土に埋まった部分を
\VUgras{礎} \textsuperscript{(2)} という。礎に囲まれた場所が\VUgras{穴倉}
\textsuperscript{(3)} をなす。この穴倉には\VUgras{風抜き窓} \textsuperscript{(5)} と
いう小さな窓があり、\VUgras{戸} \textsuperscript{(4)} と階段がある。

%%K t05-01-13 p
\VUgras{左官} \textsuperscript{(108)} が壁を積み上げていった。\VUgras{戸}
\textsuperscript{(8)} と\VUgras{窓} \textsuperscript{(17)} のための穴を残しながら。

%%K t05-01-14 p
\textsc{ジャン}。--- でもその左官が見えません。

%%K t05-01-15 p
\textsc{ルブラン氏}。--- 家の仕事はもう終わったのだ。あそこの\VUgras{馬屋}
\textsuperscript{(54)} のところにいる。自分の\VUgras{足場} \textsuperscript{(110)} に
のぼって、\VUgras{洗濯場} \textsuperscript{(56)} の壁を積んでいる。

%%K t05-01-16 p
あの男は鏝と練り舟と\VUgras{下げ振} \textsuperscript{(109)} を持っている。もっとも
この名前はお前には覚えられまい。

%%K t05-01-17 p
\textsc{ジャン}。--- 覚えられます。叔父さんに小さな鏝と練り舟をねだりますし、
下げ振は紐で自分で作りますから。

%%K t05-tit-3 sub
\VUsaut{3.0mm}
\VUpk{10.2pt}{40}{材料。}
\VUsaut{3.0mm}

%%K t05-01-18 p
\textsc{叔父}。--- なるほど結構。だが言ってごらん、ジャン、家を建てるには何が
要る。

%%K t05-01-19 p
\textsc{ジャン}。--- \VUgras{切石} \textsuperscript{(59)} と\VUgras{漆喰}
\textsuperscript{(65)} が要ります。

%%K t05-01-20 p
\textsc{叔父}。--- そのとおり。あの大きな帽子の男を見てごらん、自分の
\VUgras{荷車} \textsuperscript{(136)} に乗っている。あれは\VUgras{車力}
\textsuperscript{(135)} だ。毎日ここへ\VUgras{石の塊} \textsuperscript{(60)} を運んで
くる。荷車を中庭で下ろすと、\VUgras{石工} \textsuperscript{(83)} が塊を刻み、
自分たちの\VUgras{石鋸} \textsuperscript{(85)} で挽く。\VUgras{手伝}
\textsuperscript{(146)} がそれを左官のところへ運び、左官がそれを据える。

%%K t05-01-21 p
そのあいだ\VUgras{漆喰練り} \textsuperscript{(87)} が漆喰をこしらえる。
\VUgras{砂} \textsuperscript{(62)}、\VUgras{石灰} \textsuperscript{(63)}、
\VUgras{水} \textsuperscript{(64)} が要る。石灰は\VUgras{袋} \textsuperscript{(71)}
に入ってそばにある。\VUgras{左官の手伝} \textsuperscript{(111)} が\VUgras{手押車}
\textsuperscript{(112)} で砂を運んでき、\VUgras{弟子} \textsuperscript{(113)} は
ポンプ \textsuperscript{(58)} のところにいる。\VUgras{手桶} \textsuperscript{(114)}
を一つ満たしている。漆喰はじきにできる。\VUgras{漆喰運び}
\textsuperscript{(107)} がそれを取って、梯子で足場まで担ぎ上げる。そこでは一人の
左官が家のその面を仕上げている。

%%K t05-01-22 p
さてそれでは、\VUgras{屋根} \textsuperscript{(31)} を作る職人は何というかね。

%%K t05-01-23 p
\textsc{ジャン}。--- \VUgras{屋根葺} \textsuperscript{(118)} です。

%%K t05-tit-4 sub
\VUsaut{3.0mm}
\VUpk{10.2pt}{40}{家の屋根。}
\VUsaut{3.0mm}

%%K t05-01-24 p
\textsc{ルブラン氏}。--- そのとおり。だがまず\VUgras{大工} \textsuperscript{(115)}
が来る。

%%K t05-01-25 p
大工は\VUgras{小屋組} \textsuperscript{(35)} を作る。\VUgras{梁}
\textsuperscript{(37)} を\VUgras{込栓} \textsuperscript{(117)} でつなぐ。上にいる
幅広の天鵞絨のズボンをはいた職人が大工である。一人は小さな梁を支え、もう一人は
自分の\VUgras{斧} \textsuperscript{(116)} で込栓を削っている。\VUgras{煙突}
\textsuperscript{(41)} のそばにいるのが屋根葺である。\VUgras{根太}(*) の上に
桟を打ち、桟の上に\VUgras{スレート} \textsuperscript{(66)} を打つ。ときには瓦を
葺くこともある。

%%K t05-noto-1 noto
\VUnotes{143.2mm}{%
(*) \VUgras{Balk-o} = ドイツ語 \textit{Balken}、英語 \textit{joist}、
フランス語 \textit{solive}、イタリア語 \textit{travicello}、ロシア語
\textit{balka}、スペイン語とポルトガル語 \textit{viga}。まだ採用されていない
技術の語根であるが、フランス語 \textit{solive} の意を伝えるためにここに提案
する。それは \textit{trabo}（フランス語 \textit{poutre}）でもなく、小さな梁で
もない。床と天井とを支える角材である。%
}

%%K t05-01-26 p
馬屋の屋根は\VUgras{瓦} \textsuperscript{(55)} であり、家の屋根は
\VUgras{スレート} \textsuperscript{(32)} であり、涼み廊の屋根は\VUgras{亜鉛}
\textsuperscript{(24)} である。

%%K t05-01-27 p
\textsc{ジャン}。--- 亜鉛の屋根も屋根葺が作るのですか。

%%K t05-01-28 p
\textsc{ルブラン氏}。--- いや、それは\VUgras{亜鉛職} \textsuperscript{(121)} か
\VUgras{管職}の仕事だ。いま家の屋根に\VUgras{屋根窓} \textsuperscript{(38)} を
はめているあの男だよ。\VUgras{軒樋} \textsuperscript{(43)} と\VUgras{竪樋}
\textsuperscript{(42)} をつけるのもあの男だ。亜鉛も鉄板も鑞づけする。
\VUgras{避雷針} \textsuperscript{(40)} と\VUgras{風見} \textsuperscript{(39)} を
\VUgras{切妻} \textsuperscript{(36)} に据えたのもあの男である。

%%K t05-01-29 p
\textsc{ジャン}。--- では家はできあがったのですか。

%%K t05-tit-5 sub
\VUsaut{3.0mm}
\VUpk{10.2pt}{40}{道具。}
\VUsaut{3.0mm}

%%K t05-01-30 p
\textsc{ルブラン氏}。--- まだ少し残っている。\VUgras{壁塗} \textsuperscript{(99)}
が\VUgras{煉瓦} \textsuperscript{(61)} で仕切壁を積み、家を一間一間に分ける。
\VUgras{漆喰} \textsuperscript{(70)} を\VUgras{天井} \textsuperscript{(26)} と壁に
塗る。いまは\VUgras{涼み廊} \textsuperscript{(33)} の天井をやっている。左官と同じく
\VUgras{鏝} \textsuperscript{(100)} と\VUgras{練り舟} \textsuperscript{(101)} を
持っている。

%%K t05-01-31 p
それから指物師が戸と窓と\VUgras{板の床} \textsuperscript{(16)} と\VUgras{鎧戸}
\textsuperscript{(19)} を作る。

%%K t05-01-32 p
いまは中庭で自分の\VUgras{仕事台} \textsuperscript{(90)} で働いている。
\VUgras{鋸} \textsuperscript{(94)} --- \VUgras{木} \textsuperscript{(69)} を
挽くためのもの ---、\VUgras{鉋} \textsuperscript{(91)}、\VUgras{締具} \textsuperscript{(92)}、
\VUgras{鑿} \textsuperscript{(94 \textit{bis})}、\VUgras{コンパス}
\textsuperscript{(95)}、それに木の\VUgras{槌} \textsuperscript{(95 \textit{bis})} を
持っている。

%%K t05-01-33 p
\textsc{ジャン}。--- おや、指物師は左官よりもっと面白そうです。叔父さん、
仕事台と鋸と鉋を買ってくださいますか。メドールに小屋を作ってやりたいのです。

%%K t05-01-34 p
\textsc{叔父}。--- よろしい、坊や。

%%K t05-01-35 p
\textsc{ジャン}。--- うれしいなあ。ところで表の戸のそばにいるあの人は誰ですか。

%%K t05-tit-6 sub
\VUsaut{3.0mm}
\VUpk{10.2pt}{40}{塗り。}
\VUsaut{3.0mm}

%%K t05-01-36 p
\textsc{ルブラン氏}。--- あれは\VUgras{塗師} \textsuperscript{(102)} だ。梯子に
のぼって、\VUgras{塗料} \textsuperscript{(103)} の缶と自分の\VUgras{刷毛}
\textsuperscript{(104)} を持っている。表の戸の木に塗料を塗って、長持ちさせるので
ある。

%%K t05-01-37 p
部屋の壁紙を貼るのもあの男である。

%%K t05-01-38 p
\VUgras{表具師} \textsuperscript{(107)} は\VUgras{梯子} \textsuperscript{(106)} に
のぼって\VUgras{露台} \textsuperscript{(28)} のところにおり、\VUgras{日除}
\textsuperscript{(29)} を取りつけている。

%%K t05-01-39 p
大きな窓のそばにいる職人は\VUgras{硝子屋} \textsuperscript{(96)} である。
\VUgras{硝子} \textsuperscript{(18)} を\VUgras{パテ} \textsuperscript{(73)} ではめて
いる。あのもう一人、穴倉の戸のそばに膝をついている職人は\VUgras{錠前屋}
\textsuperscript{(97)} である。\VUgras{錠} \textsuperscript{(6)} を取りつけている。
彼の\VUgras{鑢} \textsuperscript{(98)} がそばに置いてある。

%%K t05-01-40 p
\textsc{ジャン}。--- 自分の\VUgras{金槌} \textsuperscript{(84)} で鉄を打っている
あの人も錠前屋ですか。

%%K t05-01-41 p
\textsc{ルブラン氏}。--- 正しく言えばあれは\VUgras{鍛冶} \textsuperscript{(125)}
だ。\VUgras{鉄物} \textsuperscript{(67)} を鍛えている。\VUgras{電気屋}
\textsuperscript{(124)} のためで、その電気屋は\VUgras{車庫} \textsuperscript{(51)} の
屋根の上にいる。鍛冶は\VUgras{炉} \textsuperscript{(126)}、\VUgras{金床}
\textsuperscript{(127)}、\VUgras{やっとこ} \textsuperscript{(128)} を持っている。

%%K t05-01-42 p
電気屋は電話の\VUgras{銅線} \textsuperscript{(44)} を張っている。

%%K t05-tit-7 sub
\VUpk{10.2pt}{40}{庭仕事。}
\VUsaut{3.0mm}

%%K t05-01-43 p
\textsc{叔父}。--- \VUgras{中庭} \textsuperscript{(45)} の奥、\VUgras{鉄柵}
\textsuperscript{(47)} と\VUgras{門} \textsuperscript{(48)} のそばに\VUgras{庭師}
\textsuperscript{(129)} が見えるだろう。あの小さな\VUgras{庭} \textsuperscript{(49)}
をこしらえている。もう自分の\VUgras{鍬} \textsuperscript{(131)} で土を起こし、
種を蒔き、\VUgras{熊手} \textsuperscript{(130)} でならしている。あとで自分の
\VUgras{如露} \textsuperscript{(132)} で\VUgras{花壇} \textsuperscript{(150)} に水を
やれば、草が伸びるであろう。

%%K t05-01-44 p
馬の世話と餌やりをすませたばかりの\VUgras{馬丁} \textsuperscript{(133)} が
\VUgras{馬車} \textsuperscript{(52)} を洗っている。海綿と\VUgras{手押ポンプ}
\textsuperscript{(134)} と水桶を使っている。あとでそれを車庫に入れ、重い
\VUgras{閂} \textsuperscript{(53)} で戸を閉めるであろう。

%%K t05-01-45 p
さあ、これで満足したろう。説明はもう十分だ。

%%K t05-01-46 p
\textsc{ジャン}。--- はい叔父さん。でも家の中も見たいのです。

%%K t05-01-47 p
\textsc{叔父}。--- それは明日にしよう。ルー氏が図面を説き、部屋ごとの名を教えて
くださる。

%%K t05-tit-8 sub
\VUsaut{3.0mm}
\VUpk{10.2pt}{40}{家の中の間取り。}
\VUsaut{2.4mm}

%%K t05-apar-2 apar
{\centering\textit{（図面を見よ。）}\par}
\VUsaut{2.4mm}

%%K t05-01-48 p
\textsc{建築家}。--- おいでジャン、家のいろいろな部分を見せてあげよう。

%%K t05-01-49 p
まず\VUgras{一階} \textsuperscript{(7)} に入ろう。これが\VUgras{呼鈴}の押釦
\textsuperscript{(11)} だ。三つの\VUgras{段} \textsuperscript{(14)} をのぼる。
\VUgras{階前} \textsuperscript{(9)} の段である。\VUgras{敷居} \textsuperscript{(10)}
を跨ぐ。\VUgras{表の戸} \textsuperscript{(8)} の敷居である。すると\VUgras{階段}
\textsuperscript{(13)} の下に出る。

%%K t05-01-50 p
\textsc{ジャン}。--- これが廊下ですね。

%%K t05-01-51 p
\textsc{建築家}。--- いや、これは\VUgras{玄関の間} \textsuperscript{(12)} だ。
\VUgras{廊下} \textsuperscript{(a)} は奥にある。右手にこれが\VUgras{客間}
\textsuperscript{(b)} と\VUgras{食堂} \textsuperscript{(c)}。食堂の向いが
\VUgras{台所} \textsuperscript{(d)} と\VUgras{配膳室} \textsuperscript{(e)}。その先が
\VUgras{書斎} \textsuperscript{(f)} である。\VUgras{涼み廊} \textsuperscript{(23)} は
客間の隣で、\VUgras{硝子戸} \textsuperscript{(25)} がついている。

%%K t05-01-52 p
さて階段をのぼろう。\VUgras{手すり} \textsuperscript{(15)} につかまって段を数え
ごらん。十四段ある。これが\VUgras{踊り場} \textsuperscript{(m)}。我々は
\VUgras{二階} \textsuperscript{(27)} にいる。

%%K t05-tit-9 sub
\VUsaut{3.0mm}
\VUpk{10.2pt}{40}{二階。}
\VUsaut{3.0mm}

%%K t05-01-53 p
階段の向いが\VUgras{便所} \textsuperscript{(20)} と\VUgras{化粧室}
\textsuperscript{(j)} である。右手に立派な\VUgras{寝室} \textsuperscript{(g)} が
あり、二つの窓が家の\VUgras{正面} \textsuperscript{(1)} に向いている。左手が
\VUgras{湯殿} \textsuperscript{(1)} と\VUgras{客用の間} \textsuperscript{(i)} で、
そこには大きな\VUgras{寝床の凹み} \textsuperscript{(h)} がある。寝室の隣が
\VUgras{子供部屋} \textsuperscript{(k)} である。広い\VUgras{露台}
\textsuperscript{(21)} に通じている。その露台の下にもう一つ便所
\textsuperscript{(20)} があり、壁に添えてこれが食事を知らせる\VUgras{鐘}
\textsuperscript{(22)} である。

%%K t05-01-54 p
\textsc{ジャン}。--- \VUgras{三階}へ行きましょう。

%%K t05-01-55 p
\textsc{建築家}。--- 三階はない。屋根の下には\VUgras{屋根裏部屋}
\textsuperscript{(33)} と物置 \textsuperscript{(34)} しかない。ひとまわりすんだ。
下りてよろしい。

%%K t05-01-56 p
\textsc{ジャン}。--- ありがとうございます。たいそう面白うございました。見たこと
をみな父に話します。
\VUsaut{4.0mm}
\VUfilet{22mm}
